\section{Choosing a system for a workload}
\label{sec:guidance}

The preceding sections establish the tradeoffs structurally
(\Cref{sec:axes}) and empirically (\Cref{sec:results}). This section
turns them into deployment guidance. The three axes are
\emph{hash-agnostic} --- they concern the arithmetization, the proximity
test, and the reduction field, not SHA-256 specifically --- so the
guidance applies across the bitwise-hash design space and to
heterogeneous statements that pair a hash with other arithmetic. We
assume throughout that the adder-carry soundness gap of
\Cref{sec:cmp-soundness} has been closed (e.g.\ by the lookup adder),
so the two systems are compared at equal soundness.

\subsection{The governing question}
\label{sec:guid-question}

Three properties of the workload drive the choice, each mapping to one
axis of \Cref{tab:axes}: (i) \emph{what dominates} the statement --- a
bitwise hash, or heavy big-integer / field arithmetic; (ii) the
\emph{instance scale} --- moderate versus large $\nu$; and (iii) the
\emph{output constraints} --- proof size, verifier cost, on-chain
verification. The headline summary is that hash-dominated work at scale
favours \zinc{}, a hard small-proof requirement favours \binius{} (or
recursion over \zinc{}), heterogeneous statements favour \zinc{}'s
multi-ring composition, and moderate instances are close enough that the
output constraint decides.

\subsection{Hash-dominated workloads at scale (e.g.\ hashing megabytes)}
\label{sec:guid-hash}

This is the regime where \zinc{}'s advantages compound, and three of the
four axes point the same way. \emph{Memory} is usually the binding
constraint: hashing megabytes is large-$\nu$ proving, where \binius{}'s
$\K$-valued sumcheck tensors grow super-linearly (\Cref{sec:results}
records $\sim$$27$\,GB at $\sim$$1$\,MB of SHA on a $16$\,GB host), while
\zinc{}'s base-field, bit-packed discharge is $\sim$$2\times$ leaner and
sub-linear --- so \binius{} becomes infeasible on commodity hardware
sooner. \emph{Prover throughput} favours \zinc{} for the same structural
reason it wins the discharge by $\sim$$2\times$ (\Cref{sec:res-phases}):
when the statement is just hashing, the discharge \emph{is} the cost, and
\binius{} spends roughly half its prover on the shift reduction that
\zinc{} obtains as a free $\psi$ read-off (\Cref{sec:shifts}). \emph{Verifier}
cost favours \zinc{}'s $\sqrt N$ proximity test over \binius{}'s $O(N)$
structure check, which reaches seconds at $\nu=20$.

The one \binius{} advantage, the polylog proof, flips the decision only
under a hard small-proof constraint, and even then a recursion layer over
\zinc{} is usually preferable (\Cref{sec:guid-proof}). At \emph{true}
multi-megabyte scale neither system holds the trace in memory; one chunks
the input and recurses (a hash tree is natural here), and \zinc{} is the
better per-chunk building block on both prover time and footprint.

\paragraph{Which hash, within $\Ftwo[X]$.} The hash should be chosen for
its discharge cost and soundness. \textbf{Keccak} is the natural fit:
being purely bitwise, it has \emph{no modular additions}, hence no carry
chains, hence no adder Hadamards --- its only nonlinearity is the $\chi$
AND step. It is therefore (a) fully sound with no trusted-adder caveat
at all (the gap of \Cref{sec:cmp-soundness} is specific to SHA's carries,
which Keccak lacks); (b) the leanest discharge (only $\chi$); and (c) a
clean fit for the base-field fast lane on its $64$-bit lanes. SHA-256 is
the choice only when interoperability pins it; Blake3's tree structure is
attractive when parallelism across very large inputs is decisive, at the
cost of being the most addition-heavy of the three.

\subsection{Small-proof or on-chain-verified workloads}
\label{sec:guid-proof}

Here \binius{}'s recursive proximity gives a polylog proof natively,
the one place its commitment axis dominates. But its $O(N)$ verifier
(\Cref{sec:shifts,sec:results}) makes \emph{raw} \binius{} unsuitable for
on-chain verification without its own recursion or a holographic
structure commitment (\Cref{sec:ana-air}). The practical recipe for a
small \emph{and} cheaply-verified proof is a recursion/wrapping layer
over either system, and \zinc{}'s $\sqrt N$ verifier is the cheaper inner
argument to wrap than \binius{}'s $O(N)$ one. The distinction to draw for
a client is therefore between a small \emph{standalone} proof (favours
\binius{}) and a small \emph{verified} proof (favours
\zinc{}-plus-recursion).

\subsection{Heterogeneous statements: a hash plus other arithmetic}
\label{sec:guid-hetero}

Many deployed statements pair a hash with non-bitwise arithmetic, the
archetype being signature verification: a message is hashed and the
digest is bound into a modular-exponentiation (RSA) or elliptic-curve
(ECDSA) check. The two systems answer this differently, and the
difference is itself a comparison point. \binius{}'s word circuit
expresses everything as one general circuit --- paying the $O(N)$
verifier for that generality --- whereas \zinc{}'s polynomial-ring
framing lets each sub-computation live in its \emph{natural ring},
composed in a single constraint system.

The guidance is to \emph{unify the arithmetic domain to the dominant
sub-computation rather than forcing one field}. A bitwise hash wants
characteristic two ($\Ftwo[X]$); big-integer modular arithmetic wants the
integer/limb domain. For a signature-dominated statement the proven
pattern keeps everything in the limb domain: \zinc{} already composes an
integer-domain SHA-256 with an ECDSA check in one constraint system over
limb cells $\mathrm{Int}\langle\ell\rangle$, with modular reduction
discharged by the same ideal check used elsewhere. RSA verification is
the natural extension --- a wider limb cell, square-and-multiply for the
small public exponent, and reduction modulo the public modulus by the
ideal check --- and is arguably simpler than the elliptic-curve case, as
it needs no point arithmetic. The fast $\Ftwo[X]$ hash earns its
cross-domain bridge only when the \emph{hash}, not the signature,
dominates the statement; otherwise the unified integer domain is simpler
and avoids a bit-decomposition bridge between rings.

\subsection{Summary}
\label{sec:guid-summary}

\begin{table}[h]
\centering
\renewcommand{\arraystretch}{1.3}
\begin{tabularx}{\textwidth}{@{}llX@{}}
\toprule
\textbf{Workload} & \textbf{Recommend} & \textbf{Why} \\
\midrule
Hash-dominated, large $\nu$, off-chain verify &
  \zinc{} &
  wins prover (discharge $\sim$$2\times$), memory ($\sim$$2\times$,
  sub-linear), and verifier ($\sqrt N$ vs $O(N)$) \\
Hash-dominated, small / on-chain proof &
  \zinc{} $+$ recursion &
  polylog proof needed; \zinc{}'s $\sqrt N$ verifier is cheaper to wrap
  than \binius{}'s $O(N)$. Raw \binius{} only if recursion is impossible \\
Heterogeneous (hash $+$ big-integer / field, e.g.\ RSA, ECDSA) &
  \zinc{}, unified limb domain &
  multi-ring composition; modular reduction by the ideal check; hash and
  signature in one constraint system \\
Moderate $\nu$, general &
  output-constraint decides &
  prover is a near-tie; choose \binius{} if proof size binds, \zinc{} if
  verifier cost binds \\
\bottomrule
\end{tabularx}
\caption{Workload-driven choice between the two systems, assuming equal
soundness. The recurring pattern: \zinc{} wins where prover throughput,
memory, or verifier cost binds (and at scale); \binius{} wins where a
small standalone proof binds.}
\label{tab:guidance}
\end{table}

The pattern across \Cref{tab:guidance} is the one the whole comparison
predicts: \zinc{} is the better choice wherever prover throughput,
memory, or verifier cost is the binding constraint --- which includes the
large-scale and heterogeneous regimes --- and \binius{} is the better
choice wherever a small standalone proof is the binding constraint and
recursion is not on the table.
